\documentclass[10pt,a4paper]{article}
	
	\usepackage{amsmath,amssymb,amsthm,mathtools,amsfonts}
	\usepackage{breqn}
	\usepackage[utf8]{inputenc}
	\usepackage[english]{babel}
	\usepackage[T1]{fontenc}
	\usepackage{graphicx}
	\graphicspath{{./img/}}
	
	\usepackage{fontspec}
	\setmainfont[Numbers=OldStyle]{EB Garamond}
	\newfontfamily\plantinb[Numbers=OldStyle]{Plantin MT Pro Bold Condensed}
	\newfontfamily\plantin[Numbers=OldStyle]{Plantin MT Pro Semibold}
	
	\usepackage[pdfborder={0 0 0}]{hyperref}
	\usepackage{multicol}
	\usepackage{tabularx}
	\usepackage{enumitem}
	\usepackage{wrapfig}
		\setlength{\columnsep}{0pt}
		\setlength{\intextsep}{-10pt}
	\usepackage{caption}
	\usepackage[svgnames]{xcolor}
	\usepackage{dirtytalk}
	\usepackage[modulo]{lineno}
	\usepackage{tasks}
	\newcommand\svar[1]{\newline\textcolor{red}{\textbf{Svar}\\#1}}
	\newcommand{\qm}[1]{‘‘#1”}
	\newcommand{\sqm}[1]{‘#1’}
	
	\definecolor{hgblue}{RGB}{10, 40, 80}
	\definecolor{hgred}{RGB}{140, 10, 10}
	\definecolor{hggrey}{RGB}{215, 215, 210}
	\definecolor{hggreen}{RGB}{145, 205, 205}
	
	\usepackage{titlesec}
	\titleformat{\subsection}{\color{hgred}\large}{}{}{}
	\usepackage{ulem}
	
	\setlength{\parskip}{0.5\baselineskip}
	\setlength{\parindent}{0cm}
	\setlist[enumerate]{label=\alph*),left=20pt}
	
	\usepackage{titling}
	\setlength{\droptitle}{-12ex}
	\pretitle{\begin{flushleft}\plantinb\huge\color{hgblue}}
	\posttitle{\par\end{flushleft}}
	\preauthor{\begin{flushleft}\Large}
	\postauthor{\end{flushleft}}
	\predate{\begin{flushleft}}
	\postdate{\end{flushleft}}
	
	\setcounter{secnumdepth}{0}

\begin{document}
\title{20-Minute Exerciser: \textit{Victory}}
\date{2015}
\author{Sherman Alexie}
\maketitle
    
\thispagestyle{plain}
\setcounter{page}{1}
\linenumbers

\textbf{Goal:} Analyse how guilt, poverty, pain, and religion are presented in the poem.

\textbf{Time:} 20 minutes

\subsection{1. Quick start (3 min)}
Write 3--4 lines:
\begin{enumerate}
    \item What is the speaker's situation in the poem?
    \item Why does the title \textit{Victory} seem surprising after a first reading?
\end{enumerate}

\subsection{2. Close reading (6 min)}
Answer briefly in full sentences.
\begin{enumerate}
    \item What does the stolen pair of basketball shoes represent for the speaker?
    \item Look at the repetition in \qm{guilt, guilt, guilt}. What effect does this have?
    \item Explain the effect of the exclamations: \qm{O, immoral shoes!} and \qm{O, torn skin! O, bloody heels and toes!}
    \item Find one example of imagery connected to pain. What does it add to the poem?
\end{enumerate}

\subsection{3. Deeper interpretation (7 min)}
Choose \textbf{two} of the following and answer in 4--6 lines each.
\begin{enumerate}
    \item How is poverty shown as more than just a lack of money?
    \item What kind of father is presented in the final part of the poem?
    \item Why does the poem mix religious language with basketball language?
    \item In what sense can the ending be read as a \textit{victory}?
\end{enumerate}

\subsection{4. Exit ticket (4 min)}
Write one short analytical paragraph (6--8 lines):

\qm{In \textit{Victory}, Sherman Alexie shows that true victory is not about winning in sport, but about \underline{\hspace{4cm}}. He presents this through \underline{\hspace{4cm}} and \underline{\hspace{4cm}}.}

\vspace{1em}
\textbf{Challenge if you finish early:} Find one contrast in the poem (for example hope/pain, guilt/forgiveness, sport/religion) and explain how it shapes the message.

\end{document}
